\label{kap:analyticka}
\mathquote{Each problem that I solved became a rule which served afterwards to solve other problems.}{René Descartes}

Analytická geometrie je místo, kde se setkává geometrie s algebrou. Body popíšeme čísly (souřadnicemi), přímky a kružnice rovnicemi a geometrické úlohy pak řešíme počítáním. Místo kružítka a pravítka používáme soustavy rovnic z kapitol~\ref{kap:rovnice} a~\ref{kap:kvadraticke}. Tahle myšlenka je tak silná, že na ní dnes stojí počítačová grafika, navigace i fyzika. Váš telefon totiž neumí rýsovat. Umí jen počítat.

\section{Soustava souřadnic}

\begin{definitionbox}
  \textbf{Kartézská soustava souřadnic} v rovině je tvořena dvěma navzájem kolmými číselnými osami $x$ a $y$ se společným počátkem $O$ a stejnou jednotkou. Každému bodu $A$ roviny odpovídá právě jedna uspořádaná dvojice $[a_1, a_2]$ jeho \textbf{souřadnic}. Rovina tak odpovídá kartézskému součinu $\mathbb{R}\times\mathbb{R}$ z kapitoly 2.

  \vspace{0.2cm}

  \textbf{Vzdálenost bodů} $A = [a_1, a_2]$ a $B = [b_1, b_2]$:
  \begin{equation*}
    |AB| = \sqrt{{(b_1 - a_1)}^2 + {(b_2 - a_2)}^2}.
  \end{equation*}
  \textbf{Střed úsečky} $AB$:
  \begin{equation*}
    S_{AB} = \left[\frac{a_1 + b_1}{2}, \frac{a_2 + b_2}{2}\right].
  \end{equation*}
\end{definitionbox}

Vzorec pro vzdálenost je jen Pythagorova věta. Úsečka $AB$ je přeponou pravoúhlého trojúhelníku, jehož odvěsny jsou rovnoběžné s osami a mají délky $|b_1 - a_1|$ a $|b_2 - a_2|$.
Představte si, že jdete po městě s pravoúhlými ulicemi z $A$ do $B$. Po ulicích ujdete $3$ bloky doprava a $4$ nahoru, ale pták letí přímo, po přeponě, a uletí jen $5$. Střed úsečky je zase jen `průměr': jeho souřadnice $x$ leží přesně v půli mezi $a_1$ a $b_1$, souřadnice $y$ v půli mezi $a_2$ a $b_2$.

\begin{center}
  \begin{tikzpicture}[scale=0.65, >=stealth]
    \draw[gray!30] (-0.5,-0.5) grid (5.5,6.5);
    \draw[->] (-0.5,0) -- (5.8,0) node[right] {\small $x$};
    \draw[->] (0,-0.5) -- (0,6.8) node[above] {\small $y$};
    \foreach \x in {1,...,5} \node[below] at (\x,0) {\scriptsize $\x$};
    \foreach \y in {1,...,6} \node[left] at (0,\y) {\scriptsize $\y$};
    \fill[colorofwisdom!15] (1,2) -- (4,2) -- (4,6) -- cycle;
    \draw[thick] (1,2) -- (4,2) node[midway, below] {\scriptsize $|b_1 - a_1| = 3$};
    \draw[thick] (4,2) -- (4,6) node[midway, right] {\scriptsize $|b_2 - a_2| = 4$};
    \draw (3.7,2) -- (3.7,2.3) -- (4,2.3);
    \draw[very thick, colorofwisdom] (1,2) -- (4,6) node[pos=0.8, left] {\small $5$};
    \draw[dotted, thick] (2.5,4) -- (2.5,0) (2.5,4) -- (0,4);
    \node[below] at (2.5,-0.3) {\scriptsize $\frac{5}{2}$};
    \filldraw (1,2) circle (2pt) node[below left] {\small $A$};
    \filldraw (4,6) circle (2pt) node[above right] {\small $B$};
    \filldraw[colorofwisdom] (2.5,4) circle (2pt) node[above left=-1pt] {\small $S_{AB}$};
  \end{tikzpicture}
\end{center}

\begin{example}
  Pro $A = [1, 2]$ a $B = [4, 6]$ je $|AB| = \sqrt{9 + 16} = 5$ a $S_{AB} = \left[\frac{5}{2}, 4\right]$.
\end{example}

\section{Vektory}

\begin{definitionbox}
  \textbf{Vektor} $\vec{u} = \overrightarrow{AB}$ je určen orientovanou úsečkou s počátečním bodem $A$ a koncovým bodem $B$. Dvě orientované úsečky určují stejný vektor, pokud mají stejnou délku, jsou rovnoběžné a stejně orientované. Vektor je tedy `posunutí', nezáleží na tom, odkud ho kreslíme.

  \vspace{0.2cm}

  \textbf{Souřadnice vektoru} $\overrightarrow{AB}$ jsou
  \begin{equation*}
    \vec{u} = B - A = (b_1 - a_1, b_2 - a_2).
  \end{equation*}
  \textbf{Velikost vektoru} $\vec{u} = (u_1, u_2)$ je $|\vec{u}| = \sqrt{u_1^2 + u_2^2}$. \textbf{Nulový vektor} $\vec{o} = (0, 0)$ má velikost $0$. \textbf{Jednotkový vektor} má velikost $1$.
\end{definitionbox}

Vektor si představte jako pokyn na mapě pokladu: `tři kroky na východ a jeden na sever'. Pokyn sám neříká, kde stojíte. Ať vyrazíte odkudkoli, posunete se stejně. Všechny šipky na obrázku jsou proto \textbf{tentýž} vektor $\vec{u} = (3, 1)$, jen nakreslený z různých míst.

\begin{center}
  \begin{tikzpicture}[scale=0.65, >=stealth]
    \draw[gray!30] (-0.5,-0.5) grid (9.5,4.5);
    \foreach \a/\b in {1/1, 0/3, 5/2.5, 6/0} \draw[->, very thick, colorofwisdom] (\a,\b) -- ++(3,1);
    \draw[dashed] (1,1) -- (4,1) node[midway, below] {\scriptsize $3$};
    \draw[dashed] (4,1) -- (4,2) node[midway, right] {\scriptsize $1$};
    \filldraw (1,1) circle (2pt) node[below left] {\small $A$};
    \filldraw (4,2) circle (2pt) node[above] {\small $B$};
    \node[above, colorofwisdom] at (1.5,3.5) {\small $\vec{u}$};
    \node[above, colorofwisdom] at (6.5,3) {\small $\vec{u}$};
    \node[above, colorofwisdom] at (7.5,0.5) {\small $\vec{u}$};
  \end{tikzpicture}
\end{center}

Souřadnice bodů píšeme do hranatých závorek, souřadnice vektorů do kulatých. Je to drobnost, ale zkoušející si na ní rádi zakládají, takže si na ni zvykněte.

\begin{definitionbox}
  \textbf{Operace s vektory}: Pro $\vec{u} = (u_1, u_2)$, $\vec{v} = (v_1, v_2)$ a $k\in\mathbb{R}$:
  \begin{itemize}
    \item \textbf{součet}: $\vec{u} + \vec{v} = (u_1 + v_1, u_2 + v_2)$. Geometricky: vektor $\vec{v}$ připojíme na konec $\vec{u}$ (nebo úhlopříčka rovnoběžníku),
    \item \textbf{násobek}: $k\vec{u} = (ku_1, ku_2)$. Geometricky: prodloužení $|k|$-krát, pro $k < 0$ i s otočením směru,
    \item \textbf{opačný vektor}: $-\vec{u} = (-u_1, -u_2)$.
  \end{itemize}
  Vektory $\vec{u}, \vec{v}$ jsou \textbf{kolineární} (lineárně závislé), pokud je jeden násobkem druhého. Nenulové vektory $(u_1, u_2)$ a $(v_1, v_2)$ jsou kolineární, právě když
  \begin{equation*}
    \begin{vmatrix} u_1 & v_1\\ u_2 & v_2 \end{vmatrix} = u_1v_2 - u_2v_1 = 0.
  \end{equation*}
\end{definitionbox}

\begin{center}
  \begin{tikzpicture}[scale=0.9, >=stealth]
    \coordinate (O) at (0,0); \coordinate (U) at (3,1); \coordinate (V) at (1,2);
    \draw[->, thick, colorofwisdom] (O) -- (U) node[midway, below right] {\small $\vec{u}$};
    \draw[->, thick] (O) -- (V) node[midway, left] {\small $\vec{v}$};
    \draw[->, dashed, colorofwisdom] (V) -- ($(U)+(V)$);
    \draw[->, dashed] (U) -- ($(U)+(V)$);
    \draw[->, very thick] (O) -- ($(U)+(V)$) node[above right] {\small $\vec{u} + \vec{v}$};
  \end{tikzpicture}
\end{center}

Násobek vektoru číslem si představte jako změnu rychlosti při stejném směru: $2\vec{u}$ je dvakrát delší krok, $\frac{1}{2}\vec{u}$ poloviční krok a $-\vec{u}$ krok stejně dlouhý, ale pozpátku.

\begin{center}
  \begin{tikzpicture}[scale=0.7, >=stealth]
    \draw[gray!30] (-0.5,-0.5) grid (11.5,2.5);
    \draw[->, very thick, colorofwisdom] (0,0) -- (2,1) node[above] {\small $\vec{u}$};
    \draw[->, very thick] (2.5,0) -- (6.5,2) node[above] {\small $2\vec{u}$};
    \draw[->, very thick] (7,0) -- (8,0.5) node[above] {\small $\frac{1}{2}\vec{u}$};
    \draw[->, very thick] (11,1) -- (9,0) node[below] {\small $-\vec{u}$};
  \end{tikzpicture}
\end{center}

Ano, ten determinant je ten samý jako v kapitole~\ref{kap:rovnice}. Tam jsme si řekli, že jeho absolutní hodnota je obsah rovnoběžníku, který vektory vytvoří. Kolineární vektory vytvoří `rovnoběžník' s nulovým obsahem, a proto je determinant nulový.

\begin{center}
  \begin{tikzpicture}[scale=0.8, >=stealth]
    \draw[gray!30] (-0.5,-0.5) grid (4.5,3.5);
    \fill[colorofwisdom!20] (0,0) -- (3,1) -- (4,3) -- (1,2) -- cycle;
    \draw[dashed] (3,1) -- (4,3) -- (1,2);
    \draw[->, very thick, colorofwisdom] (0,0) -- (3,1) node[below right] {\small $\vec{u} = (3, 1)$};
    \draw[->, very thick] (0,0) -- (1,2) node[above left] {\small $\vec{v} = (1, 2)$};
    \node at (2,1.5) {\small $S = 5$};
    \node at (2,-1) {\scriptsize $\begin{vmatrix} 3 & 1\\ 1 & 2 \end{vmatrix} = 6 - 1 = 5$};
  \end{tikzpicture}
  \hspace{1.5cm}
  \begin{tikzpicture}[scale=0.8, >=stealth]
    \draw[gray!30] (-0.5,-0.5) grid (4.5,3.5);
    \draw[->, very thick] (0,0) -- (4,2) node[above] {\small $\vec{v} = (4, 2)$};
    \draw[->, very thick, colorofwisdom] (0,0) -- (2,1) node[pos=0.6, below right=2pt] {\small $\vec{u} = (2, 1)$};
    \node at (2,-1) {\scriptsize $\begin{vmatrix} 2 & 4\\ 1 & 2 \end{vmatrix} = 4 - 4 = 0$};
  \end{tikzpicture}
\end{center}
Vpravo jsou vektory kolineární a rovnoběžník se `splácl' do úsečky.

\subsection{Skalární součin}

\begin{definitionbox}
  \textbf{Skalární součin} vektorů $\vec{u} = (u_1, u_2)$ a $\vec{v} = (v_1, v_2)$ je \textbf{číslo}
  \begin{equation*}
    \vec{u}\cdot\vec{v} = u_1v_1 + u_2v_2.
  \end{equation*}
  \textbf{Věta}: Pro úhel $\varphi\in[0^\circ, 180^\circ]$ mezi nenulovými vektory platí
  \begin{equation*}
    \vec{u}\cdot\vec{v} = |\vec{u}|\cdot|\vec{v}|\cdot\cos\varphi,\qquad\text{tedy}\qquad\cos\varphi = \frac{\vec{u}\cdot\vec{v}}{|\vec{u}|\cdot|\vec{v}|}.
  \end{equation*}
  Speciálně: nenulové vektory jsou \textbf{kolmé}, právě když $\vec{u}\cdot\vec{v} = 0$. Dále $\vec{u}\cdot\vec{u} = |\vec{u}|^2$.
\end{definitionbox}

Co skalární součin `měří'? Představte si, že vektor $\vec{v}$ je tyč a slunce svítí kolmo na přímku, na které leží $\vec{u}$. Stín tyče na této přímce má délku $|\vec{v}|\cos\varphi$, a skalární součin je tedy $|\vec{u}|$ krát délka stínu. Říká, \textit{jak moc míří $\vec{v}$ stejným směrem jako $\vec{u}$}. Je-li úhel tupý, stín padá dozadu a součin je záporný. Je-li pravý, stín je jen bod a součin je nula. Odtud ta nejdůležitější vlastnost: kolmost poznáme podle nulového skalárního součinu.

\begin{center}
  \begin{tikzpicture}[scale=0.75, >=stealth]
    \foreach \ang/\sx/\lab in {50/0/{$\vec{u}\cdot\vec{v} > 0$}, 90/5.2/{$\vec{u}\cdot\vec{v} = 0$}, 130/10.4/{$\vec{u}\cdot\vec{v} < 0$}} {
      \begin{scope}[xshift=\sx cm]
        \draw[gray] (-1.5,0) -- (3.3,0);
        \draw[dashed, gray] (\ang:2) -- ({2*cos(\ang)},0);
        \draw[line width=2.5pt, colorofwisdom!60] (0,0) -- ({2*cos(\ang)},0);
        \draw[->, very thick, colorofwisdom] (0,0) -- (3,0) node[below] {\small $\vec{u}$};
        \draw[->, very thick] (0,0) -- (\ang:2) node[above] {\small $\vec{v}$};
        \draw (0,0) ++(0:0.45) arc[start angle=0, end angle=\ang, radius=0.45];
        \node at ({\ang/2}:0.75) {\scriptsize $\varphi$};
        \node at (0.7,-0.9) {\small \lab};
      \end{scope}
    }
  \end{tikzpicture}
\end{center}
Tlustý úsek na přímce vektoru $\vec{u}$ je `stín' vektoru $\vec{v}$, tedy jeho kolmý průmět.

\begin{proof}
  Vektory $\vec{u}$, $\vec{v}$ umístíme do společného počátku. Třetí strana trojúhelníku, který vytvoří, je vektor $\vec{u} - \vec{v}$.
  \begin{center}
    \begin{tikzpicture}[scale=0.8, >=stealth]
      \coordinate (O) at (0,0); \coordinate (U) at (4,0.5); \coordinate (V) at (1,2);
      \fill[colorofwisdom!10] (O) -- (U) -- (V) -- cycle;
      \draw[->, very thick, colorofwisdom] (O) -- (U) node[midway, below] {\small $\vec{u}$};
      \draw[->, very thick] (O) -- (V) node[midway, left] {\small $\vec{v}$};
      \draw[->, very thick, gray] (V) -- (U) node[midway, above right] {\small $\vec{u} - \vec{v}$};
      \draw (O) ++(7.1:0.6) arc[start angle=7.1, end angle=63.4, radius=0.6];
      \node at (35:0.9) {\scriptsize $\varphi$};
    \end{tikzpicture}
  \end{center}
  Šedý vektor $\vec{u} - \vec{v}$ vede z konce $\vec{v}$ do konce $\vec{u}$, protože $\vec{v} + (\vec{u} - \vec{v}) = \vec{u}$. Podle kosinové věty z kapitoly~\ref{kap:goniometrie}:
  \begin{equation*}
    |\vec{u} - \vec{v}|^2 = |\vec{u}|^2 + |\vec{v}|^2 - 2|\vec{u}||\vec{v}|\cos\varphi.
  \end{equation*}
  Zároveň v souřadnicích:
  \begin{equation*}
    |\vec{u} - \vec{v}|^2 = {(u_1 - v_1)}^2 + {(u_2 - v_2)}^2 = |\vec{u}|^2 + |\vec{v}|^2 - 2(u_1v_1 + u_2v_2).
  \end{equation*}
  Porovnáním dostaneme $u_1v_1 + u_2v_2 = |\vec{u}||\vec{v}|\cos\varphi$.
\end{proof}

\begin{example}
  Vektory $\vec{u} = (3, 4)$ a $\vec{v} = (4, -3)$ mají skalární součin $12 - 12 = 0$, jsou tedy kolmé. Mimochodem, vektor kolmý k $(a, b)$ dostaneme vždy tak, že souřadnice prohodíme a jedné změníme znaménko: $(-b, a)$ nebo $(b, -a)$.

  Úhel mezi $\vec{u} = (1, 2)$ a $\vec{v} = (3, 1)$: $\cos\varphi = \frac{3 + 2}{\sqrt{5}\cdot\sqrt{10}} = \frac{5}{\sqrt{50}} = \frac{\sqrt{2}}{2}$, tedy $\varphi = 45^\circ$.
\end{example}

\begin{problem}
  \begin{enumerate}[label=(\alph*)]
    \item Pro $A = [-1, 3]$, $B = [2, -1]$ určete $\overrightarrow{AB}$, $|AB|$ a střed úsečky $AB$.
    \item Pro která $k$ jsou vektory $(2, k)$ a $(k, 8)$ kolineární? Pro které $k$ je $(2, k)$ kolmý na $(3, -6)$?
    \item Určete vnitřní úhly trojúhelníku $A = [0, 0]$, $B = [4, 0]$, $C = [1, 3]$.
    \item Doplňte bod $D$ tak, aby $ABCD$ byl rovnoběžník, kde $A = [1, 1]$, $B = [5, 2]$, $C = [6, 5]$.
  \end{enumerate}
\end{problem}

\begin{solution}
  \begin{enumerate}[label=(\alph*)]
    \item $\overrightarrow{AB} = (3, -4)$, $|AB| = 5$, $S = \left[\frac{1}{2}, 1\right]$.
    \item $2\cdot 8 - k\cdot k = 0$, $k = \pm 4$. Kolmost: $6 - 6k = 0$, $k = 1$.
    \item U $A$: $\overrightarrow{AB} = (4, 0)$, $\overrightarrow{AC} = (1, 3)$, $\cos\alpha = \frac{4}{4\sqrt{10}} = \frac{1}{\sqrt{10}}$, $\alpha\approx 71{,}6^\circ$. U $B$: $\overrightarrow{BA} = (-4, 0)$, $\overrightarrow{BC} = (-3, 3)$, $\cos\beta = \frac{12}{4\cdot 3\sqrt{2}} = \frac{\sqrt{2}}{2}$, $\beta = 45^\circ$. Pak $\gamma\approx 63{,}4^\circ$.
    \item $D = A + \overrightarrow{BC} = [1, 1] + (1, 3) = [2, 4]$.
  \end{enumerate}
  \begin{center}
    \begin{tikzpicture}[scale=0.6, >=stealth]
      \draw[gray!30] (-0.5,-0.5) grid (6.5,5.5);
      \fill[colorofwisdom!10] (1,1) -- (5,2) -- (6,5) -- (2,4) -- cycle;
      \draw[thick] (1,1) -- (5,2) (6,5) -- (2,4);
      \draw[->, very thick, colorofwisdom] (5,2) -- (6,5) node[midway, right] {\scriptsize $\overrightarrow{BC}$};
      \draw[->, very thick, colorofwisdom] (1,1) -- (2,4) node[midway, left] {\scriptsize $\overrightarrow{BC}$};
      \foreach \p/\c/\pos in {A/{1,1}/below left, B/{5,2}/below right, C/{6,5}/above right, D/{2,4}/above left} \filldraw (\c) circle (2pt) node[\pos] {\small $\p$};
    \end{tikzpicture}
  \end{center}
  Pozor na pořadí vrcholů: rovnoběžník $ABCD$ se obchází dokola, takže $D$ leží naproti $B$, ne naproti $C$.
\end{solution}

\section{Přímka v rovině}

Přímku můžeme popsat třemi způsoby. Každý se hodí na něco jiného, a proto je potřeba umět mezi nimi přecházet.

\begin{definitionbox}
  \textbf{Parametrické vyjádření}: Přímka procházející bodem $A$ se \textbf{směrovým vektorem} $\vec{u}\neq\vec{o}$ je množina bodů
  \begin{equation*}
    X = A + t\vec{u},\quad t\in\mathbb{R},\qquad\text{tedy}\qquad x = a_1 + tu_1,\quad y = a_2 + tu_2.
  \end{equation*}

  \textbf{Obecná rovnice}: Přímka s \textbf{normálovým vektorem} $\vec{n} = (a, b)\neq\vec{o}$ (kolmým na přímku) má rovnici
  \begin{equation*}
    ax + by + c = 0.
  \end{equation*}

  \textbf{Směrnicový tvar}: Pokud přímka není rovnoběžná s osou $y$, lze ji zapsat jako
  \begin{equation*}
    y = kx + q,
  \end{equation*}
  kde $k = \tg\varphi$ je \textbf{směrnice} ($\varphi$ je úhel s kladnou poloosou $x$) a $q$ je úsek na ose $y$. To je přesně graf lineární funkce z kapitoly~\ref{kap:elementarni}.
\end{definitionbox}

Parametrické vyjádření si můžete představit jako pohyb. V čase $t = 0$ jste v bodě $A$ a každou sekundu se posunete o vektor $\vec{u}$. Za všechny časy (i záporné, tedy i v minulosti) opíšete celou přímku. Obecná rovnice naopak přímku nepopisuje jako cestu, ale jako podmínku: bod leží na přímce, právě když splňuje rovnici.

\begin{example}
  Najděme všechna vyjádření přímky $AB$, kde $A = [1, 2]$ a $B = [4, 6]$.

  Směrový vektor $\vec{u} = (3, 4)$. Parametricky: $x = 1 + 3t$, $y = 2 + 4t$.

  Normálový vektor je kolmý na $\vec{u}$: $\vec{n} = (4, -3)$. Obecná rovnice $4x - 3y + c = 0$, dosazením bodu $A$: $4 - 6 + c = 0$, tedy $c = 2$ a $4x - 3y + 2 = 0$.

  Směrnicový tvar: $y = \frac{4}{3}x + \frac{2}{3}$.

  Z parametrického vyjádření dostaneme obecnou rovnici také tak, že z obou rovnic vyloučíme $t$: $t = \frac{x - 1}{3} = \frac{y - 2}{4}$, tedy $4x - 4 = 3y - 6$ a $4x - 3y + 2 = 0$.
\end{example}

Na obrázku je přímka z příkladu. Body s parametrem $t = 0$ a $t = 1$ jsou $A$ a $B$, polovina parametru dává střed úsečky $AB$ a záporné $t$ leží `za' bodem $A$. Směrový vektor $\vec{u}$ míří podél přímky, normálový vektor $\vec{n}$ je na ni kolmý. Obecná rovnice tedy říká: \textit{bod $X$ leží na přímce, právě když je vektor $\overrightarrow{AX}$ kolmý na $\vec{n}$}, tedy $\vec{n}\cdot(X - A) = 0$. Po roznásobení je to přesně $4x - 3y + 2 = 0$.

\begin{center}
  \begin{tikzpicture}[scale=0.55, >=stealth]
    \draw[gray!30] (-1.5,-1.5) grid (6.5,8.5);
    \draw[->] (-1.5,0) -- (6.8,0) node[right] {\small $x$};
    \draw[->] (0,-1.5) -- (0,8.8) node[above] {\small $y$};
    \draw[thick] (-1.25,-1) -- (6.25,9) node[below right] {\small $p$};
    \draw[->, very thick, colorofwisdom] (1,2) -- (4,6) node[pos=0.3, right=4pt] {\small $\vec{u} = (3, 4)$};
    \draw[->, very thick] (1,2) -- (5,-1) node[right] {\small $\vec{n} = (4, -3)$};
    \draw (1.24,2.32) -- (1.56,2.08) -- (1.32,1.76);
    \foreach \t/\lab in {-0.5/{-\frac{1}{2}}, 0/0, 0.5/{\frac{1}{2}}, 1/1, 1.5/{\frac{3}{2}}} {
      \filldraw[colorofwisdom] ({1+3*\t},{2+4*\t}) circle (2.5pt) node[left, xshift=-2pt, fill=white, inner sep=1pt] {\scriptsize $t = \lab$};
    }
    \node[above left=-1pt, xshift=-2pt, yshift=5pt] at (1,2) {\small $A$};
    \node[above left=-1pt, xshift=-2pt, yshift=5pt] at (4,6) {\small $B$};
  \end{tikzpicture}
\end{center}

\begin{definitionbox}
  \textbf{Vzájemná poloha dvou přímek}: Pro přímky $p: a_1x + b_1y + c_1 = 0$ a $q: a_2x + b_2y + c_2 = 0$:
  \begin{itemize}
    \item jsou \textbf{různoběžné}, pokud normálové vektory nejsou kolineární. Průsečík najdeme vyřešením soustavy dvou rovnic,
    \item jsou \textbf{rovnoběžné různé}, pokud jsou normálové vektory kolineární, ale rovnice nejsou násobkem jedna druhé,
    \item jsou \textbf{totožné}, pokud je jedna rovnice násobkem druhé.
  \end{itemize}
  To je přesně rozbor počtu řešení soustavy dvou lineárních rovnic z kapitoly~\ref{kap:rovnice}. Tam jsme řekli, že soustava popisuje dvě přímky, a teď to víme jistě.

  \vspace{0.2cm}

  \textbf{Odchylka přímek} $\varphi\in[0^\circ, 90^\circ]$:
  \begin{equation*}
    \cos\varphi = \frac{|\vec{n}_1\cdot\vec{n}_2|}{|\vec{n}_1|\cdot|\vec{n}_2|}\qquad\text{(nebo totéž se směrovými vektory)}.
  \end{equation*}
  Absolutní hodnota zajistí, že dostaneme ostrý úhel. Přímky v směrnicovém tvaru jsou rovnoběžné, právě když $k_1 = k_2$, a kolmé, právě když $k_1\cdot k_2 = -1$.

  \vspace{0.2cm}

  \textbf{Vzdálenost bodu} $M = [m_1, m_2]$ od přímky $ax + by + c = 0$:
  \begin{equation*}
    d = \frac{|am_1 + bm_2 + c|}{\sqrt{a^2 + b^2}}.
  \end{equation*}
\end{definitionbox}

\begin{center}
  \begin{tikzpicture}[scale=0.8]
    \draw[thick, colorofwisdom] (0,0) -- (3,2);
    \draw[thick] (0,2) -- (3,0.5);
    \filldraw (1.875,1.25) circle (2pt);
    \node at (1.5,-0.5) {\scriptsize různoběžné};
    \begin{scope}[xshift=4cm]
      \draw[thick, colorofwisdom] (0,0.3) -- (3,1.8);
      \draw[thick] (0,1) -- (3,2.5);
      \node at (1.5,-0.5) {\scriptsize rovnoběžné různé};
    \end{scope}
    \begin{scope}[xshift=8cm]
      \draw[line width=3pt, colorofwisdom!50] (0,0.3) -- (3,1.8);
      \draw[thick] (0,0.3) -- (3,1.8);
      \node at (1.5,-0.5) {\scriptsize totožné};
    \end{scope}
  \end{tikzpicture}
\end{center}

U odchylky můžeme klidně použít normálové vektory místo směrových. Normálový vektor je směrový vektor otočený o $90^\circ$, a když obě přímky otočíme o stejný pravý úhel, úhel mezi nimi se nezmění. Jediný háček je, že vektory mohou svírat tupý úhel $180^\circ - \varphi$ (stačí jeden z nich otočit opačným směrem). Proto je v čitateli absolutní hodnota.

\begin{center}
  \begin{tikzpicture}[scale=1, >=stealth]
    \draw[thick, colorofwisdom] (200:2) -- (20:2.6) node[right] {\small $p_1$};
    \draw[thick] (255:1.2) -- (75:2.2) node[above] {\small $p_2$};
    \draw (20:0.6) arc[start angle=20, end angle=75, radius=0.6];
    \node at (47.5:0.85) {\scriptsize $\varphi$};
    \draw[->, very thick, colorofwisdom] (0,0) -- (110:1.6) node[above] {\small $\vec{n}_1$};
    \draw[->, very thick] (0,0) -- (165:1.6) node[left] {\small $\vec{n}_2$};
    \draw (110:0.9) arc[start angle=110, end angle=165, radius=0.9];
    \node at (137.5:1.15) {\scriptsize $\varphi$};
    \draw[gray] (20:0.25) -- ($(20:0.25)+(110:0.25)$) -- (110:0.25);
  \end{tikzpicture}
\end{center}

Vzdálenost bodu od přímky je délka \textbf{nejkratší} cesty k ní, a ta vede kolmo, tedy ve směru normálového vektoru. Představte si, že stojíte na louce v bodě $M$ a přímka je řeka. K řece nejrychleji doběhnete kolmo k břehu, do bodu $Q$. Důkaz níže tu kolmou cestu nemusí vůbec hledat. Stačí vzít libovolný bod $P$ na břehu a z vektoru $\overrightarrow{PM}$ si nechat jen tu část, která míří ve směru $\vec{n}$ (jeho `stín' na normále, viz skalární součin). Na obrázku je to situace z příkladu níže: $M = [3, 1]$, přímka $4x - 3y + 2 = 0$ a $P = \left[\frac{5}{2}, 4\right]$. Body $P$, $P'$, $M$, $Q$ tvoří obdélník, takže $|PP'| = |QM| = d$.

\begin{center}
  \begin{tikzpicture}[scale=1.2, >=stealth]
    \draw[gray!30] (0,0.5) grid (5,4.5);
    \fill[colorofwisdom!12] (2.5,4) -- (4.26,2.68) -- (3,1) -- (1.24,2.32) -- cycle;
    \draw[thick] (0.25,1) -- (2.875,4.5) node[right] {\small $4x - 3y + 2 = 0$};
    \draw[->, thick, gray] (2.5,4) -- (3,1) node[pos=0.25, left] {\scriptsize $\overrightarrow{PM}$};
    \draw[very thick, colorofwisdom] (3,1) -- (1.24,2.32) node[midway, below left] {\small $d$};
    \draw[->, very thick, colorofwisdom!60] (2.5,4) -- (4.26,2.68) node[midway, above right] {\small $\vec{n}$};
    \draw[dashed] (4.26,2.68) -- (3,1);
    \draw ($(1.24,2.32)+(0.15,0.2)$) -- ++(0.2,-0.15) -- ++(-0.15,-0.2);
    \draw ($(4.26,2.68)+(-0.2,0.15)$) -- ++(-0.15,-0.2) -- ++(0.2,-0.15);
    \filldraw (3,1) circle (1.6pt) node[below right] {\small $M = [3, 1]$};
    \filldraw (2.5,4) circle (1.6pt) node[above left] {\small $P$};
    \filldraw (1.24,2.32) circle (1.6pt) node[left] {\small $Q$};
    \filldraw (4.26,2.68) circle (1.6pt) node[right] {\small $P'$};
  \end{tikzpicture}
\end{center}

\begin{proof}[(vzorce pro vzdálenost)]
  Nechť $P = [p_1, p_2]$ je libovolný bod přímky, tedy $ap_1 + bp_2 = -c$. Vzdálenost bodu $M$ od přímky je délka průmětu vektoru $\overrightarrow{PM}$ do směru kolmého na přímku, tedy do směru normálového vektoru $\vec{n}$. Z definice skalárního součinu je tato délka
  \begin{equation*}
    d = \frac{|\overrightarrow{PM}\cdot\vec{n}|}{|\vec{n}|} = \frac{|a(m_1 - p_1) + b(m_2 - p_2)|}{\sqrt{a^2 + b^2}} = \frac{|am_1 + bm_2 - (ap_1 + bp_2)|}{\sqrt{a^2 + b^2}} = \frac{|am_1 + bm_2 + c|}{\sqrt{a^2 + b^2}}.
  \end{equation*}
\end{proof}

\begin{proof}[(podmínky kolmosti $k_1k_2 = -1$)]
  Přímka $y = kx + q$ má směrový vektor $(1, k)$. Dvě přímky jsou kolmé, právě když jsou kolmé jejich směrové vektory, tedy $(1, k_1)\cdot(1, k_2) = 1 + k_1k_2 = 0$.
\end{proof}

\begin{example}
  \begin{itemize}
    \item Průsečík přímek $x + y - 3 = 0$ a $2x - y = 0$: sečtením $3x = 3$, tedy $x = 1$ a $y = 2$. Průsečík je $[1, 2]$.
    \item Vzdálenost bodu $M = [3, 1]$ od přímky $4x - 3y + 2 = 0$ je $\frac{|12 - 3 + 2|}{\sqrt{16 + 9}} = \frac{11}{5}$.
    \item Odchylka přímek $x - y = 0$ a $y = 0$: $\vec{n}_1 = (1, -1)$, $\vec{n}_2 = (0, 1)$, $\cos\varphi = \frac{1}{\sqrt{2}}$, $\varphi = 45^\circ$.
  \end{itemize}
\end{example}

\begin{definitionbox}
  \textbf{Polorovina}: Přímka $ax + by + c = 0$ rozdělí rovinu na dvě poloroviny. Body, pro které $ax + by + c > 0$, leží v jedné z nich, body s $ax + by + c < 0$ ve druhé. Normálový vektor $\vec{n} = (a, b)$ míří do té, kde je výraz kladný.
\end{definitionbox}

\begin{center}
  \begin{tikzpicture}[scale=0.7, >=stealth]
    \fill[colorofwisdom!12] (-0.5,4.5) -- (4.5,4.5) -- (4.5,-0.5) -- cycle;
    \draw[gray!30] (-0.5,-0.5) grid (4.5,4.5);
    \draw[->] (-0.5,0) -- (4.8,0) node[right] {\small $x$};
    \draw[->] (0,-0.5) -- (0,4.8) node[above] {\small $y$};
    \draw[thick] (-0.5,4.5) -- (4.5,-0.5) node[right] {\small $x + y - 4 = 0$};
    \draw[->, very thick, colorofwisdom] (1.5,2.5) -- (2.5,3.5) node[above] {\small $\vec{n} = (1, 1)$};
    \filldraw (0,0) circle (2pt) node[below left] {\scriptsize $[0, 0]$};
    \filldraw (3,3) circle (2pt) node[right] {\scriptsize $[3, 3]$};
    \node[colorofwisdom] at (3.9,1.5) {\scriptsize $> 0$};
    \node at (0.9,1) {\scriptsize $< 0$};
  \end{tikzpicture}
\end{center}
Polorovinu si představte jako jednu stranu plotu. Dosazením bodu do levé strany rovnice zjistíme, na které straně plotu stojí. Normálový vektor ukazuje `ven' na kladnou stranu.

\begin{problem}
  \begin{enumerate}[label=(\alph*)]
    \item Napište obecnou rovnici přímky, která prochází bodem $P = [2, -1]$ a je (i) rovnoběžná, (ii) kolmá k přímce $3x + 2y - 5 = 0$.
    \item Určete vzájemnou polohu přímek $p: x = 1 + 2t$, $y = 3 - t$ a $q: x + 2y - 7 = 0$.
    \item Určete vzdálenost rovnoběžek $3x - 4y + 1 = 0$ a $3x - 4y - 9 = 0$.
    \item Určete odchylku přímek $y = 2x + 1$ a $y = -3x + 4$.
    \item Leží body $[0, 0]$ a $[3, 3]$ ve stejné polorovině vzhledem k přímce $x + y - 4 = 0$?
  \end{enumerate}
\end{problem}

\begin{solution}
  \begin{enumerate}[label=(\alph*)]
    \item (i) $3x + 2y + c = 0$, $6 - 2 + c = 0$, tedy $3x + 2y - 4 = 0$. (ii) Normálový vektor kolmice je směrovým vektorem původní přímky, $(2, -3)$: $2x - 3y + c = 0$, $4 + 3 + c = 0$, tedy $2x - 3y - 7 = 0$.
    \item Směrový vektor $p$ je $(2, -1)$, normálový tedy $(1, 2)$, stejný jako u $q$. Bod $[1, 3]$ přímky $p$ dosadíme do $q$: $1 + 6 - 7 = 0$. Leží na ní, takže přímky jsou \textbf{totožné}.
    \item Vezmeme bod první přímky, například $[1, 1]$ ($3 - 4 + 1 = 0$), a spočítáme jeho vzdálenost od druhé: $\frac{|3 - 4 - 9|}{5} = 2$. Obecně pro rovnoběžky $ax + by + c_1 = 0$ a $ax + by + c_2 = 0$ je vzdálenost $\frac{|c_1 - c_2|}{\sqrt{a^2 + b^2}}$.
    \item Směrové vektory $(1, 2)$ a $(1, -3)$, $\cos\varphi = \frac{|1 - 6|}{\sqrt{5}\sqrt{10}} = \frac{5}{\sqrt{50}} = \frac{\sqrt{2}}{2}$, $\varphi = 45^\circ$.
    \item Dosadíme: $0 + 0 - 4 = -4 < 0$ a $3 + 3 - 4 = 2 > 0$. Body leží v opačných polorovinách.
  \end{enumerate}
\end{solution}

\section{Aplikace na trojúhelník}

\begin{definitionbox}
  Pro trojúhelník $ABC$ se souřadnicemi vrcholů:
  \begin{itemize}
    \item \textbf{těžiště}: $T = \left[\frac{a_1 + b_1 + c_1}{3}, \frac{a_2 + b_2 + c_2}{3}\right]$,
    \item \textbf{obsah}: $S = \frac{1}{2}\left|\det\left(\overrightarrow{AB}, \overrightarrow{AC}\right)\right| = \frac{1}{2}|(b_1 - a_1)(c_2 - a_2) - (b_2 - a_2)(c_1 - a_1)|$.
  \end{itemize}
\end{definitionbox}

Vzorec pro těžiště plyne z toho, že těžiště dělí těžnici v poměru $2:1$. Vzorec pro obsah z toho, že trojúhelník je polovina rovnoběžníku určeného vektory $\overrightarrow{AB}$ a $\overrightarrow{AC}$.

\begin{center}
  \begin{tikzpicture}[scale=0.85, >=stealth]
    \fill[colorofwisdom!10] (0,0) -- (3.5,0.5) -- (4.5,2.5) -- (1,2) -- cycle;
    \fill[colorofwisdom!30] (0,0) -- (3.5,0.5) -- (1,2) -- cycle;
    \draw[dashed] (3.5,0.5) -- (4.5,2.5) -- (1,2);
    \draw[thick] (3.5,0.5) -- (1,2);
    \draw[->, very thick] (0,0) -- (3.5,0.5) node[below right] {\small $B$};
    \draw[->, very thick] (0,0) -- (1,2) node[above left] {\small $C$};
    \node[below left] at (0,0) {\small $A$};
    \node at (1.5,0.85) {\small $S$};
    \node at (3,1.65) {\small $S$};
  \end{tikzpicture}
  \hspace{1.5cm}
  \begin{tikzpicture}[scale=0.85]
    \coordinate (A) at (0,0); \coordinate (B) at (4,0); \coordinate (C) at (1.2,2.7);
    \draw[thick] (A) -- (B) -- (C) -- cycle;
    \draw[gray] (A) -- ($(B)!0.5!(C)$) (B) -- ($(A)!0.5!(C)$);
    \draw[very thick, colorofwisdom] (C) -- ($(C)!0.6667!($(A)!0.5!(B)$)$) node[midway, right] {\scriptsize $2$};
    \draw[very thick, colorofwisdom!50] ($(C)!0.6667!($(A)!0.5!(B)$)$) -- ($(A)!0.5!(B)$) node[midway, right] {\scriptsize $1$};
    \filldraw ($(A)!0.5!(B)$) circle (1.5pt);
    \filldraw[colorofwisdom] (1.7333,0.9) circle (2pt) node[left=2pt] {\small $T$};
    \node[below left] at (A) {\small $A$}; \node[below right] at (B) {\small $B$}; \node[above] at (C) {\small $C$};
  \end{tikzpicture}
\end{center}
Vlevo: rovnoběžník je rozdělen úhlopříčkou na dva shodné trojúhelníky, takže trojúhelník má polovinu obsahu. Vpravo: těžiště $T$ dělí každou těžnici v poměru $2:1$, tedy leží ve třetině těžnice od strany. Představte si trojúhelník z kartonu: v těžišti ho udržíte na špičce tužky.

\begin{example}
  Trojúhelník $A = [0, 0]$, $B = [4, 0]$, $C = [1, 3]$ má těžiště $T = \left[\frac{5}{3}, 1\right]$ a obsah $S = \frac{1}{2}|4\cdot 3 - 0\cdot 1| = 6$. Kontrola: základna $AB$ má délku $4$ a výška z $C$ je $3$, tedy $S = \frac{4\cdot 3}{2} = 6$.

  Rovnice výšky z vrcholu $C$: je kolmá na $AB$, takže její normálový vektor je $\overrightarrow{AB} = (4, 0)$, a rovnice je $4x + c = 0$, po dosazení $C$: $x = 1$. Rovnice výšky z $A$: normálový vektor je $\overrightarrow{BC} = (-3, 3)$, neboli $(1, -1)$, rovnice $x - y = 0$. Průsečík obou výšek je ortocentrum $V = [1, 1]$.
\end{example}

\begin{center}
  \begin{tikzpicture}[scale=1.1]
    \draw[gray!30] (-0.5,-0.5) grid (4.5,3.5);
    \draw[->] (-0.5,0) -- (4.8,0) node[right] {\small $x$};
    \draw[->] (0,-0.5) -- (0,3.8) node[above] {\small $y$};
    \fill[colorofwisdom!10] (0,0) -- (4,0) -- (1,3) -- cycle;
    \draw[very thick] (0,0) -- (4,0) -- (1,3) -- cycle;
    \draw[thick, colorofwisdom, dashed] (1,3) -- (1,0) (0,0) -- (2,2) (4,0) -- (0.4,1.2);
    \draw[gray] (1,0) -- (1,0.2) -- (1.2,0.2) -- (1.2,0);
    \draw[gray] (2,2) -- (1.86,1.86) -- (2,1.72) -- (2.14,1.86);
    \draw[gray, dotted, thick] (1,3) -- (2,0) (0,0) -- (2.5,1.5) (4,0) -- (0.5,1.5);
    \filldraw[colorofwisdom] (1,1) circle (2pt) node[above left] {\small $V$};
    \filldraw (1.6667,1) circle (2pt) node[right=2pt] {\small $T$};
    \node[below left] at (0,0) {\small $A$}; \node[below right] at (4,0) {\small $B$}; \node[above] at (1,3) {\small $C$};
  \end{tikzpicture}
\end{center}
Čárkovaně jsou výšky (všechny tři se protnou v ortocentru $V$, třetí výška z $B$ má rovnici $x + 3y - 4 = 0$ a bodem $V$ opravdu prochází), tečkovaně těžnice protínající se v těžišti $T$.

\begin{problem}
  Je dán trojúhelník $A = [1, 1]$, $B = [7, 3]$, $C = [3, 7]$.
  \begin{enumerate}[label=(\alph*)]
    \item Určete jeho těžiště a obsah.
    \item Napište rovnici těžnice z vrcholu $C$.
    \item Napište rovnici osy strany $AB$.
    \item Rozhodněte, zda je trojúhelník rovnoramenný nebo pravoúhlý.
  \end{enumerate}
\end{problem}

\begin{solution}
  \begin{enumerate}[label=(\alph*)]
    \item $T = \left[\frac{11}{3}, \frac{11}{3}\right]$. $\overrightarrow{AB} = (6, 2)$, $\overrightarrow{AC} = (2, 6)$, $S = \frac{1}{2}|36 - 4| = 16$.
    \item Střed $AB$ je $[4, 2]$, směrový vektor těžnice $(1, -5)$, normálový $(5, 1)$: $5x + y + c = 0$, $20 + 2 + c = 0$, tedy $5x + y - 22 = 0$. (Kontrola: $C$ dává $15 + 7 - 22 = 0$.)
    \item Osa prochází středem $[4, 2]$ a má normálový vektor $\overrightarrow{AB} = (6, 2)$, tedy $(3, 1)$: $3x + y - 14 = 0$.
    \item $|AB| = |AC| = \sqrt{40}$, je rovnoramenný se základnou $BC$. Skalární součiny: $\overrightarrow{AB}\cdot\overrightarrow{AC} = 24\neq 0$, $\overrightarrow{BA}\cdot\overrightarrow{BC} = (-6)(-4) + (-2)(4) = 16\neq 0$, $\overrightarrow{CA}\cdot\overrightarrow{CB} = (-2)(4) + (-6)(-4) = 16\neq 0$. Pravoúhlý není.
  \end{enumerate}
  \begin{center}
    \begin{tikzpicture}[scale=0.55]
      \draw[gray!30] (-0.5,-0.5) grid (8.5,8.5);
      \draw[->] (-0.5,0) -- (8.8,0) node[right] {\small $x$};
      \draw[->] (0,-0.5) -- (0,8.8) node[above] {\small $y$};
      \fill[colorofwisdom!10] (1,1) -- (7,3) -- (3,7) -- cycle;
      \draw[very thick] (1,1) -- (7,3) -- (3,7) -- cycle;
      \draw[thick, colorofwisdom] (3,7) -- (4,2) node[pos=0.12, right=2pt] {\scriptsize těžnice};
      \draw[thick, dashed] (4.6667,0) -- (2,8) node[below left] {\scriptsize osa $AB$};
      \draw ($(4,2)+(0.2846,0.0949)$) -- ++(-0.0949,0.2846) -- ++(-0.2846,-0.0949);
      \filldraw (4,2) circle (2pt);
      \filldraw[colorofwisdom] (3.6667,3.6667) circle (2.5pt) node[right=2pt] {\small $T$};
      \node[below left] at (1,1) {\small $A$}; \node[right] at (7,3) {\small $B$}; \node[above right] at (3,7) {\small $C$};
      \draw[thick] (3.163,1.51) -- (3.037,1.89) (1.81,4.063) -- (2.19,3.937);
    \end{tikzpicture}
  \end{center}
  Čárky na stranách $AB$ a $AC$ naznačují, že jsou stejně dlouhé. Osa strany $AB$ prochází jejím středem $[4, 2]$ kolmo, těžnice z $C$ míří do téhož středu.
\end{solution}

\section{Příběh --- Lebka filozofa}

Souřadnice jsou pojmenované po Reném Descartovi (latinsky Cartesius), francouzském filozofovi a matematikovi, autorovi věty `Myslím, tedy jsem'. Podle legendy ho myšlenka souřadnic napadla, když jako nemocný ležel v posteli a pozoroval mouchu na stropě. Uvědomil si, že její polohu lze kdykoli popsat dvěma čísly: vzdáleností od dvou stěn. Descartes mimochodem celý život vstával pozdě, protože ráno v posteli prý myslel nejlépe. Čtenáři této knihy jistě chápou.

Právě to se mu stalo osudným. V roce 1649 ho švédská královna Kristina pozvala do Stockholmu, aby ji učil filozofii. Královna ale trvala na tom, že lekce budou v pět hodin ráno, v mrazivém zámku a uprostřed švédské zimy. Descartes po několika měsících dostal zápal plic a v únoru 1650 zemřel. Někteří historici dnes tvrdí, že ho otrávil arzenem katolický kněz, který se bál, že by Descartovy názory odradily protestantskou královnu od přestupu ke katolictví. Kristina mimochodem o několik let později opravdu konvertovala a vzdala se trůnu.

Tím ale Descartův příběh neskončil. Když po šestnácti letech převáželi jeho ostatky do Francie, hlava v rakvi chyběla. Lebka se pak víc než sto let potulovala Evropou, prodávala se a kupovala. Každý majitel na ni napsal své jméno a krátkou báseň, takže je dnes popsaná latinskými verši. Dnes ji najdete v Muzeu člověka v Paříži, zatímco zbytek těla leží v kostele Saint-Germain-des-Prés na druhém břehu Seiny. Zakladatel moderní filozofie, který tvrdil, že mysl a tělo jsou dvě oddělené věci, tak skončil přesně podle své vlastní teorie.

\section{Shrnutí a otázky k ústní zkoušce}

\begin{problem}
  Odpovězte vlastními slovy:
  \begin{enumerate}
    \item Jak se počítá vzdálenost dvou bodů a střed úsečky? Odkud se vzorec pro vzdálenost bere?
    \item Co je vektor? Jak se určí jeho souřadnice a velikost? Jak se vektory sčítají a násobí číslem?
    \item Kdy jsou dva vektory kolineární? Jak to souvisí s determinantem?
    \item Definujte skalární součin. Dokažte vztah se kosinem úhlu. Kdy jsou vektory kolmé?
    \item Jaké znáte způsoby vyjádření přímky v rovině? Co je směrový a normálový vektor? Jak se mezi vyjádřeními přechází?
    \item Jak určíte vzájemnou polohu dvou přímek a jejich průsečík?
    \item Jak se počítá odchylka dvou přímek a vzdálenost bodu od přímky? Odvoďte vzorec pro vzdálenost.
    \item Kdy jsou dvě přímky ve směrnicovém tvaru rovnoběžné a kdy kolmé?
    \item Jak spočítáte těžiště a obsah trojúhelníku ze souřadnic jeho vrcholů?
  \end{enumerate}
\end{problem}
